\section{Normalization cross-check: metric, Casimir, Laplacian, and \texorpdfstring{$b_0$}{b0}}
\label{app:normalization_crosscheck}

This appendix collects, in one place, the conventions linking:
(i) the fixed bi-invariant Riemannian metric on $G$ (equivalently an Ad-invariant inner product on $\mathfrak g$),
(ii) the Laplace--Beltrami operator $\Delta$ and quadratic Casimir eigenvalues $C_2(\pi)$,
and (iii) the quadratic coefficient $b_0$ in the small-coupling beta expansion used in the dyadic gradient-flow step scaling module.

\subsection{From the inner product to the Laplacian and the Casimir}

Fix the inner product $\langle\cdot,\cdot\rangle$ on $\mathfrak g$ from Definition~\ref{def:metric_casimir}.
Let $(X_a)_{a=1}^{d_G}$ be any $\langle\cdot,\cdot\rangle$-orthonormal basis of $\mathfrak g$, and let $\widetilde X_a$ denote
its corresponding left-invariant vector field on $G$.

\begin{lemma}[Laplacian/Casimir consistency]
\label{lem:laplacian_casimir_consistency}
The Laplace--Beltrami operator associated to the fixed bi-invariant metric is
\begin{equation}
\label{eq:laplacian_as_sum_squares}
\Delta\;=\;\sum_{a=1}^{d_G} \widetilde X_a^{\,2}.
\end{equation}
For every irreducible unitary representation $\pi$ of $G$, the Casimir element
$\Omega:=\sum_{a=1}^{d_G} X_a^2$ acts as a scalar on $\pi$:
$\,d\pi(\Omega)=-C_2(\pi)\,I$.
Equivalently, for the character $\chi_\pi$,
\begin{equation}
\label{eq:laplacian_character_eigen}
\Delta\chi_\pi\ =\ -C_2(\pi)\,\chi_\pi,
\end{equation}
which is the convention used throughout this manuscript.
\end{lemma}

\begin{proof}
Let $(X_a)_{a=1}^{d_G}$ be an orthonormal basis of $\mathfrak g$ for the fixed Ad-invariant inner product and let $\widetilde X_a$
be the corresponding left-invariant vector fields.
For a bi-invariant metric, the Levi--Civita connection satisfies
\[
\nabla_{\widetilde X}\widetilde Y=\tfrac12[\widetilde X,\widetilde Y]
\qquad\text{for all left-invariant vector fields }\widetilde X,\widetilde Y,
\]
so in particular $\nabla_{\widetilde X_a}\widetilde X_a=\tfrac12[\widetilde X_a,\widetilde X_a]=0$ for each $a$.
The Laplace--Beltrami operator on functions may be written in an orthonormal frame as
\[
\Delta f=\sum_{a=1}^{d_G}\big(\widetilde X_a(\widetilde X_a f)-(\nabla_{\widetilde X_a}\widetilde X_a)f\big),
\]
hence \eqref{eq:laplacian_as_sum_squares} follows.

Next let $\pi$ be an irreducible unitary representation and let $f_{v,w}(g):=\langle v,\pi(g)w\rangle$ be a matrix coefficient.
For $X\in\mathfrak g$,
\[
(\widetilde X f_{v,w})(g)
=\left.\frac{d}{dt}\right|_{t=0}\langle v,\pi(e^{tX}g)w\rangle
=\langle v,d\pi(X)\pi(g)w\rangle.
\]
Therefore
\[
\begin{aligned}
(\widetilde X_a^{\,2} f_{v,w})(g)
&=\langle v,d\pi(X_a)^2\,\pi(g)w\rangle,\\
(\Delta f_{v,w})(g)
&=\left\langle v,\Big(\sum_{a=1}^{d_G}d\pi(X_a)^2\Big)\pi(g)w\right\rangle\\
&=\langle v,d\pi(\Omega)\pi(g)w\rangle.
\end{aligned}
\]
By Schur's lemma the Casimir element $\Omega=\sum_a X_a^2$ acts on the irreducible $\pi$ by a scalar, and with the sign convention
of Definition~\ref{def:metric_casimir} this scalar is $-C_2(\pi)$, i.e.\ $d\pi(\Omega)=-C_2(\pi)I$.
Hence $\Delta f_{v,w}=-C_2(\pi)f_{v,w}$ for all matrix coefficients.
Summing over an orthonormal basis $v=w$ gives \eqref{eq:laplacian_character_eigen} for $\chi_\pi$.
\end{proof}

\subsection{Trace conventions and quadratic expansion of the Wilson plaquette weight}

The Wilson plaquette weight is built from the fixed Wilson representation $\rho$ from \eqref{eq:families}.
To avoid trace-normalization ambiguities, it is useful to isolate the unique constant relating the representation trace to
$\langle\cdot,\cdot\rangle$.

\begin{definition}[Dynkin index relative to the fixed metric]
\label{def:dynkin_index}
Let $\sigma$ be a finite-dimensional unitary representation of $G$ with differential $\sigma_\ast$.
The (positive) Dynkin index $I_\sigma$ relative to $\langle\cdot,\cdot\rangle$ is the unique scalar such that
\begin{equation}
\label{eq:dynkin_index_def}
-\Tr_{\sigma}(\sigma_\ast(X)\,\sigma_\ast(Y))\ =\ I_\sigma\,\langle X,Y\rangle\qquad (X,Y\in\mathfrak g).
\end{equation}
\end{definition}

\begin{lemma}[Quadratic coefficient at the identity]
\label{lem:wilson_quadratic_coefficient}
Let $\rho$ be the fixed Wilson representation used in the action.
As $X\to 0$ in $\mathfrak g$,
\begin{equation}
\label{eq:wilson_quad_expansion}
1-\frac{1}{\dim\rho}\Re\Tr_{\rho}(e^X)
\;=\;\frac{I_\rho}{2\,\dim\rho}\,\|X\|^2\ +\ O(\|X\|^3).
\end{equation}
Consequently, the Wilson potential
$V_{\beta,\alpha}(X)=\alpha\beta\bigl(1-\frac{1}{\dim\rho}\Re\Tr_{\rho}(e^X)\bigr)$ has Hessian
$\nabla^2 V_{\beta,\alpha}(0)=\alpha\beta\,(I_\rho/\dim\rho)\,I$.
\end{lemma}

\begin{proof}
Expand $e^X=I+X+\tfrac12X^2+O(\|X\|^3)$.
Since $X$ is skew-adjoint in a unitary representation, $\Re\Tr_\rho(X)=0$.
Moreover $\Re\Tr_\rho(X^2)=\Tr_\rho(X^2)$ and
$-\Tr_\rho(X^2)=I_\rho\,\|X\|^2$ by \eqref{eq:dynkin_index_def}.
Insert these into the expansion to obtain \eqref{eq:wilson_quad_expansion} and differentiate twice.
\end{proof}

\subsection{Relation between \texorpdfstring{$b_0$}{b0} and the conventional one-loop coefficient (optional)}

In the Lane~C step-scaling module, $b_0$ denotes the (scheme-invariant) quadratic coefficient in the small-$u$ expansion
\begin{equation}
\beta_{\mathrm{GF}}(u)\ =\ -2b_0\,u^2\ +\ O(u^3),
\qquad u=g^2_{\mathrm{GF}}.
\end{equation}
The load-bearing use in this submission is only the strict positivity $b_0>0$, proved internally in Appendix~F.
For orientation, we record the conventional perturbative identification.

\begin{lemma}[Matching to the conventional one-loop coefficient]
\label{lem:b0_one_loop_match}
Assume that the gradient-flow coupling satisfies
$u=g^2+O(g^4)$ as $g\to 0$ for some reference coupling $g$.
If the beta function of $g$ has the conventional one-loop form
\begin{equation}
\label{eq:beta_g_conventional}
\mu\frac{dg}{d\mu}= -\frac{b_0^{\mathrm{conv}}}{16\pi^2}\,g^3\ +\ O(g^5),
\end{equation}
then
\begin{equation}
\label{eq:b0_relation}
\beta_{\mathrm{GF}}(u)=\mu\frac{du}{d\mu}= -2\,\frac{b_0^{\mathrm{conv}}}{16\pi^2}\,u^2\ +\ O(u^3),
\end{equation}
so that $b_0=b_0^{\mathrm{conv}}/(16\pi^2)$ under the tree-level matching normalization.
In pure Yang--Mills one has $b_0^{\mathrm{conv}}=\frac{11}{3}C_2(\mathrm{ad})$ in the standard convention.
\end{lemma}

\begin{proof}
Differentiate $u=g^2+O(g^4)$ with respect to $\ln\mu$ and use \eqref{eq:beta_g_conventional}:
$\mu\,d(g^2)/d\mu = 2g\,\mu\,dg/d\mu = -2(b_0^{\mathrm{conv}}/(16\pi^2))\,g^4+O(g^6)$.
Replace $g^4$ by $u^2+O(u^3)$ to obtain \eqref{eq:b0_relation}.
\end{proof}

\begin{remark}[Scaling of numerical values]
If one rescales the inner product $\langle\cdot,\cdot\rangle\mapsto s\langle\cdot,\cdot\rangle$ with $s>0$, then the Laplacian
and all Casimir eigenvalues scale by $\Delta\mapsto s^{-1}\Delta$ and $C_2(\pi)\mapsto s^{-1}C_2(\pi)$.
A corresponding rescaling of the coupling variable restores the tree-level matching condition $u=g_0^2+O(g_0^4)$, so the sign $b_0>0$ is invariant.
The present manuscript treats $b_0$ and the $C_2$ normalization as named primitives (Appendix~\ref{app:parameter_schedule}) and uses only the resulting inequalities.
\end{remark}


\subsection{Explicit dim-$6$ plaquette / rectangle / chair coefficient cross-check}

For the audited dim-$6$ loop block in Section~\ref{sec:sharp_local}, we use the standard tree-level Symanzik / L\"uscher--Weisz matching coefficients for the normalized four-link / six-link gauge-loop basis \cite{LuscherWeisz1985,Husung2023}.
Write $(e_0,e_1,e_2,e_3)$ for the coefficients of the plaquette, rectangle, chair, and twisted-chair loops, normalized by
\begin{equation}\label{eq:appendixL_loop_norm}
e_0+8e_1+8e_2+16e_3=1.
\end{equation}
In the on-shell CP-even dim-$6$ gauge basis
\[
\mathcal G_{6}^{(1)}\sim \sum_{\mu,\nu,\rho}\Tr\big([D_\mu,F_{\nu\rho}]\,[D_\mu,F_{\nu\rho}]\big),
\qquad
\mathcal G_{6}^{(2)}\sim \sum_{\mu,\nu}\Tr\big([D_\mu,F_{\mu\nu}]\,[D_\mu,F_{\mu\nu}]\big),
\]
the tree-level coefficients are
\begin{equation}\label{eq:appendixL_loop_dim6_coeffs}
\bar c_1=\frac{e_2}{3},
\qquad
\bar c_2=e_1-e_3+\frac1{12}.
\end{equation}
For the normalized local loop densities used in the Lane~A sample sector one has
\[
P:(e_0,e_1,e_2,e_3)=(1,0,0,0),
\qquad
R:(0,\tfrac18,0,0),
\qquad
C:(0,0,\tfrac18,0),
\]
so
\begin{equation}\label{eq:appendixL_PRC_coeff_table}
P:(\bar c_1,\bar c_2)=\Bigl(0,\frac1{12}\Bigr),
\qquad
R:(\bar c_1,\bar c_2)=\Bigl(0,\frac5{24}\Bigr),
\qquad
C:(\bar c_1,\bar c_2)=\Bigl(\frac1{24},\frac1{12}\Bigr).
\end{equation}
Subtracting the plaquette contribution therefore gives the explicit improved dim-$6$ probe coefficients
\begin{equation}\label{eq:appendixL_PRC_improved_matrix}
R-P:\ (0,\tfrac18),
\qquad
C-P:\ (\tfrac1{24},0),
\end{equation}
which is the source of the concrete Lane~A sample matrix
\[
D_{6,\mathrm{tree}}^{PRC}=\begin{pmatrix}0&\tfrac18\\[0.35em]\tfrac1{24}&0\end{pmatrix},
\qquad
\det D_{6,\mathrm{tree}}^{PRC}=-\frac1{192}.
\]
This is exactly the determinant certificate used in Proposition~\ref{prop:D6_tree_certificate}.


\subsection{Packet 5--7 certification note and Lane~A interface map}
\label{subsec:LaneA_packet_interface_map}

For a first-pass audit of Packets~5--7, the only theorem-facing chain to keep in view is
\[
\begin{aligned}
\text{graded identity / block coefficients}
&\Longrightarrow \text{one-shell transport + finite-truncation inverse control}\\
&\Longrightarrow \text{finite-cap flowed-sector bridge}
\Longrightarrow \text{Lane~A closure}\\
&\Longrightarrow \text{bounded-base cyclicity}.
\end{aligned}
\]
Table~\ref{tab:LaneA_packet_interface_map} freezes that packet interface in one place: it records what each node consumes from the previous stage, what it exports to the next stage, and what it does \emph{not} yet prove. In particular, Theorem~\ref{thm:finite_cap_flowed_sector_bridge} is the first mixed-correlator closure node on the frozen spine; the promoted finite-cap bridge into Theorem~\ref{thm:LaneA_closure} is the theorem-facing state/positivity package; Theorem~\ref{thm:LaneA_closure}(b) is only the finite-cap OS-structure theorem; Corollary~\ref{cor:LaneA_bounded_state_restriction} is only the bounded-state compatibility node; the inductive-union extension does not occur until Corollary~\ref{cor:LaneA_full_scope}; and bounded-base cyclicity is added only by Theorem~\ref{thm:LaneA_bounded_base_cyclic}.

\begingroup
\small
\setlength{\tabcolsep}{4pt}
\begin{longtable}{@{}>{\RaggedRight\arraybackslash}p{0.18\textwidth}>{\RaggedRight\arraybackslash}p{0.22\textwidth}>{\RaggedRight\arraybackslash}p{0.27\textwidth}>{\RaggedRight\arraybackslash}p{0.23\textwidth}@{}}
\caption{Packet~5--7 Lane~A interface map on the frozen spine.}\label{tab:LaneA_packet_interface_map}\\
\toprule
\textbf{Node} & \textbf{Consumes} & \textbf{Exports to the next node} & \textbf{Does not yet prove}\\
\midrule
\endfirsthead
\toprule
\textbf{Node} & \textbf{Consumes} & \textbf{Exports to the next node} & \textbf{Does not yet prove}\\
\midrule
\endhead
Packet~5a: exact-dimension breadth\\
Theorem~\ref{thm:LaneA_associated_graded_identity}, Theorem~\ref{thm:LaneA_block_coefficient_theorem}
& Canonical exact-dimension quotient blocks, coherent local/flowed contraction formulas, and the renormalized field family of Theorem~\ref{thm:field_algebra_truncation}.
& Basis-free diagonal symbol control and coefficient extraction for every exact-dimension symmetry block; this is the blockwise UV datum used by the inverse-control step.
& No mixed-correlator closure, no finite-cap state, and no inductive-union sharp-local extension.\\[0.45em]

Packet~5b: one-shell / finite-truncation package\\
Theorem~\ref{thm:canonical_block_one_shell_increment}, Theorem~\ref{thm:finite_truncation_inverse_control}, Corollary~\ref{cor:finite_truncation_SFTE}
& Packet~5a block coefficients together with the insertion RG / OS bounds controlling the off-diagonal blocks.
& Canonical one-shell transport, polylogarithmic inverse control, and SFTE/remainder compatibility on every finite engineering cap; Corollary~\ref{cor:LaneA_all_finite_truncations} packages this as the unconditional finite-truncation input.
& Still no packaged mixed-cap state and no extension from $\mathcal A_{\mathrm{flow}}$ to the full sharp-local algebra.\\[0.45em]

Packet~6a: finite-cap bridge\\
Corollary~\ref{cor:finite_cap_mixed_family_packaging}, Theorem~\ref{thm:finite_cap_flowed_sector_bridge}
& Packet~5b finite-truncation inverse/remainder data together with the Lane~C flowed state and bounded positive-time base state.
& First mixed-correlator closure node: capwise packaged bounded states, inverse-flow identification, and mixed-channel Cauchy control at fixed engineering cap.
& No inductive-union passage and no bounded-base cyclicity.\\[0.45em]

Packet~6b: Lane~A closure\\
Theorem~\ref{thm:LaneA_closure}, Corollary~\ref{cor:LaneA_bounded_state_restriction}, Corollary~\ref{cor:LaneA_full_scope}
& Packet~6a finite-cap bridge together with cap-compatibility of the canonical normalization scheme.
& Sharp-local continuum state on each finite cap, then the bounded-state restriction compatibility clause, and only afterward the full inductive-union algebra.
& Does not itself prove bounded-base cyclicity.\\[0.45em]

Packet~7: bounded-base cyclicity\\
Theorem~\ref{thm:LaneA_bounded_base_cyclic}
& The full sharp-local state from Corollary~\ref{cor:LaneA_full_scope} and the common-core / graph-topology package used in the final Lane~A block.
& The original bounded positive-time base algebra remains cyclic in the reconstructed sharp-local Hilbert space.
& Adds no new inverse-flow identification and no new mixed-correlator closure statement.\\
\bottomrule
\end{longtable}
\endgroup


\subsection{Lane A canonical block basis and UV normalization cross-check}

The full-scope Lane~A theorem in Section~\ref{sec:sharp_local} uses a canonical blockwise normalization scheme rather than an ever-growing list of unrelated sector-specific choices.
This subsection records that protocol in one place so a referee can read the later coefficient theorem with a fixed dictionary.

\paragraph{Block labels.}
For each engineering dimension $\Delta$, CP parity $\epsilon$, and Euclidean tensor type $\lambda$,
write
\[
\mathscr V_{\Delta,\epsilon,\lambda}\subset \mathscr V_\Delta
\]
for the corresponding finite symmetry block from Definition~\ref{def:local_operator_spaces}.
The direct sum over all such blocks is $\mathscr V_\Delta$ itself.

\paragraph{Canonical local basis.}
Inside each block choose a basis by complete contractions of $F_{\mu\nu}$ and covariant derivatives,
modulo the quotient relations already built into Definition~\ref{def:local_operator_spaces}
(lattice symmetries, total derivatives, and Bianchi identities).
This choice is made once and for all for each exact-dimension symmetry block and is then not altered when larger engineering caps are introduced.

\begin{definition}[Canonical block norm]\label{def:LaneA_canonical_block_norm}
Fix an exact-dimension symmetry block $\mathscr V_{\Delta,\epsilon,\lambda}$ with canonical basis
$\{e_1,\dots,e_m\}$ chosen above.
For $v=\sum_{i=1}^m c_i e_i$ define
\[
\|v\|_{\mathrm{can}}:=\max_{1\le i\le m}|c_i|.
\]
If $T:\mathscr V_{\Delta,\epsilon,\lambda}\to\mathscr V_{\Delta',\epsilon',\lambda'}$ is linear, write
$\|T\|_{\mathrm{op,can}}$ for the corresponding operator norm.
Because each block is finite dimensional and its basis is fixed once and for all,
this norm is intrinsic to the chosen canonical block and does not depend on how that block is embedded into a larger truncation.
\end{definition}

\paragraph{Canonical flowed basis.}
Given that local basis, define the flowed probe basis by applying the \emph{same} contraction formulas to the flowed curvature
$G_{\mu\nu}(t,\cdot)$ and flowed covariant derivatives.

\paragraph{Filtered quotient model behind the canonical blocks.}
To make the later coefficient theorem basis-free, it is convenient to describe each exact-dimension block intrinsically rather than by coordinates.
Let $\mathscr J_{\le \Delta}$ denote the finite-dimensional real vector space of local gauge-invariant jet polynomials of engineering dimension $\le \Delta$
built from $F_{\mu\nu}$ and covariant derivatives, and let $\mathscr R_{\le \Delta}\subset \mathscr J_{\le \Delta}$ be the homogeneous subspace generated by
lattice symmetry relations, total derivatives, and Bianchi relations.
Define the exact-dimension quotient
\begin{equation}\label{eq:LaneA_exact_dim_quotient}
\mathscr Q_{\Delta}
:=
\mathscr J_{\le \Delta}/(\mathscr J_{\le \Delta-1}+\mathscr R_{\le \Delta}).
\end{equation}
Its $(\epsilon,\lambda)$-isotypic component is exactly the symmetry block $\mathscr V_{\Delta,\epsilon,\lambda}$ from Definition~\ref{def:local_operator_spaces}.
Likewise, define the flowed quotient $\mathscr Q^{\mathrm{flow}}_{\Delta}(t)$ by replacing the unflowed jets with the flowed jets $\nabla_t^{\alpha}G_{\mu\nu}(t)$ but imposing the same homogeneous quotient relations.

\begin{lemma}[Small-flow substitution respects the exact-dimension quotient]\label{lem:LaneA_flow_substitution_filtered}
Let $\Gamma_t$ be the algebra map sending each unflowed jet $\nabla^{\alpha}F_{\mu\nu}$ to the corresponding flowed jet
$\nabla_t^{\alpha}G_{\mu\nu}(t)$.
Then for every homogeneous representative $P\in \mathscr J_{\le \Delta}$ of exact engineering dimension $\Delta$ one has
\begin{equation}\label{eq:LaneA_filtered_substitution_expansion}
\Gamma_t P
=
P+\sum_{m\ge 1} t^m P_{\Delta+2m}+R_{<\Delta},
\end{equation}
where each $P_{\Delta+2m}$ is a local gauge-invariant jet polynomial of exact engineering dimension $\Delta+2m$ and
$R_{<\Delta}\in \mathscr J_{\le \Delta-1}+\mathscr R^{\mathrm{flow}}_{\le \Delta}$.
Consequently $\Gamma_t$ descends to a well-defined linear map on the exact-dimension quotient,
\[
\mathrm{gr}_{\Delta}(\Gamma_t):\mathscr Q_{\Delta}\longrightarrow \mathscr Q^{\mathrm{flow}}_{\Delta}(t).
\]
\end{lemma}

\begin{proof}
For every jet generator, the gradient-flow equation and its covariant derivative consequences give the standard small-flow expansion
\[
\nabla_t^{\alpha}G_{\mu\nu}(t)
=
\nabla^{\alpha}F_{\mu\nu}+\sum_{m\ge 1} t^m K^{(m)}_{\alpha,\mu\nu},
\]
where $K^{(m)}_{\alpha,\mu\nu}$ is a local gauge-covariant polynomial whose engineering dimension is larger by $2m$.
This is the jet-level form of the heat-kernel / Duhamel expansion already used throughout Lane~A.
Applying the same complete-contraction formula to a homogeneous monomial of dimension $\Delta$ therefore reproduces that monomial itself,
plus terms with explicit powers $t^m$ multiplying local expressions of exact dimension $\Delta+2m$.
The quotient relations are homogeneous and are respected by the substitution $F\mapsto G(t)$,
so any symmetry relation, total derivative, or Bianchi relation is sent into the corresponding flowed relation.
Thus the only surviving class in the quotient $\mathscr Q_{\Delta}$ is the original exact-dimension class of $P$,
which yields the descended map $\mathrm{gr}_{\Delta}(\Gamma_t)$.
\end{proof}

\begin{theorem}[Associated-graded identity on exact-dimension symmetry blocks]\label{thm:LaneA_associated_graded_identity}
For every exact-dimension symmetry block $\mathscr V_{\Delta,\epsilon,\lambda}\subset \mathscr Q_{\Delta}$,
the descended small-flow map of Lemma~\ref{lem:LaneA_flow_substitution_filtered} is the identity after the canonical identification of
$\mathscr Q_{\Delta}$ with $\mathscr Q^{\mathrm{flow}}_{\Delta}(t)$ by the common complete-contraction formulas:
\[
\mathrm{gr}_{\Delta}(\Gamma_t)\!\restriction_{\mathscr V_{\Delta,\epsilon,\lambda}}
=
I_{\mathscr V_{\Delta,\epsilon,\lambda}}.
\]
Equivalently, the tree-level small-flow matching on the exact-dimension quotient block is basis-free and exactly the identity.
\end{theorem}

\begin{proof}
Lemma~\ref{lem:LaneA_flow_substitution_filtered} shows that every homogeneous representative of an exact-dimension class is sent to itself,
modulo the sum of higher-dimension terms with explicit positive powers of $t$ and the homogeneous quotient relations.
Passing to the exact-dimension associated-graded projection therefore discards the higher-dimension corrections,
and the quotient \eqref{eq:LaneA_exact_dim_quotient} kills all lower-order ambiguities,
leaving precisely the original class.
Because the $(\epsilon,\lambda)$ symmetry decomposition is itself homogeneous, the same conclusion holds after restricting to the symmetry block
$\mathscr V_{\Delta,\epsilon,\lambda}$.
No basis choice is used at this stage.
\end{proof}

\begin{lemma}[Coefficient extraction is well defined on arbitrary exact-dimension quotient blocks]\label{lem:LaneA_block_coeff_extract}
Fix an exact-dimension symmetry block $\mathscr V_{\Delta,\epsilon,\lambda}\subset \mathscr Q_{\Delta}$ with canonical basis
$\{e_1,\dots,e_m\}$, and let
\[
\operatorname{Coeff}_{\Delta,\epsilon,\lambda}:\mathscr V_{\Delta,\epsilon,\lambda}\longrightarrow \mathbb R^m
\]
be the corresponding coordinate map.
Let
\[
\operatorname{Coeff}^{\mathrm{flow}}_{\Delta,\epsilon,\lambda,t}:\mathscr V^{\mathrm{flow}}_{\Delta,\epsilon,\lambda}(t)\longrightarrow \mathbb R^m
\]
denote the analogous coordinate map on the flowed block defined by the same complete-contraction formulas.
For any homogeneous representative $P\in \mathscr J_{\le \Delta}$ of exact engineering dimension $\Delta$, write
$[P]_{\Delta,\epsilon,\lambda}$ for the projection of its quotient class to $\mathscr V_{\Delta,\epsilon,\lambda}$.
Then:
\begin{enumerate}[label=(\roman*),leftmargin=2.3em]
\item $\operatorname{Coeff}_{\Delta,\epsilon,\lambda}([P]_{\Delta,\epsilon,\lambda})$ depends only on the quotient class of $P$ modulo
$\mathscr J_{\le \Delta-1}+\mathscr R_{\le \Delta}$;
\item if $P'\equiv P\pmod{\mathscr J_{\le \Delta-1}+\mathscr R_{\le \Delta}}$, then
\[
\mathrm{gr}_{\Delta}(\Gamma_t)[P']
=
\mathrm{gr}_{\Delta}(\Gamma_t)[P]
\qquad\text{in }\mathscr Q^{\mathrm{flow}}_{\Delta}(t),
\]
so coefficient extraction after small-flow substitution is likewise representative-independent;
\item under the canonical identification of unflowed and flowed exact-dimension blocks by the common complete-contraction formulas,
\begin{equation}\label{eq:LaneA_coeff_extract_identity}
\operatorname{Coeff}^{\mathrm{flow}}_{\Delta,\epsilon,\lambda,t}\Big(\big(\mathrm{gr}_{\Delta}(\Gamma_t)[P]\big)_{\Delta,\epsilon,\lambda}\Big)
=
\operatorname{Coeff}_{\Delta,\epsilon,\lambda}\big([P]_{\Delta,\epsilon,\lambda}\big).
\end{equation}
\end{enumerate}
Equivalently, lower-dimension terms and homogeneous quotient relations contribute zero to the exact-dimension coefficient vector,
both before and after the small-flow substitution.
\end{lemma}

\begin{proof}
Part~(i) is exactly the definition of the quotient \eqref{eq:LaneA_exact_dim_quotient}.
For part~(ii), write
\[
P'-P=Q+R,
\qquad
Q\in \mathscr J_{\le \Delta-1},\quad R\in \mathscr R_{\le \Delta}.
\]
Because the quotient relations are homogeneous and respected by the substitution $F\mapsto G(t)$,
Lemma~\ref{lem:LaneA_flow_substitution_filtered} shows that $\Gamma_tQ$ contributes only to lower engineering degree and that
$\Gamma_tR$ lands in the flowed relation space.
Hence $\Gamma_tP'$ and $\Gamma_tP$ define the same class in the flowed exact-dimension quotient.

For part~(iii), apply Lemma~\ref{lem:LaneA_flow_substitution_filtered} to $P$:
\[
\Gamma_tP=P+\sum_{m\ge 1} t^m P_{\Delta+2m}+R_{<\Delta}.
\]
When one passes to the exact-dimension associated-graded piece, the higher-dimension terms
$t^mP_{\Delta+2m}$ are discarded by the projection to degree $\Delta$, while $R_{<\Delta}$ vanishes in the quotient.
Thus the class of $\Gamma_tP$ in the flowed block is exactly the class determined by $P$ itself.
Applying the canonical coordinate maps gives \eqref{eq:LaneA_coeff_extract_identity}.
\end{proof}

\paragraph{Coherent renormalized basis.}
The coherent graded renormalization scheme of Theorem~\ref{thm:field_algebra_truncation} and Lemma~\ref{lem:coherent_Z}
chooses renormalized fields so that when a larger engineering cap is introduced, the fields already normalized in smaller caps are unchanged.
Equivalently, the exact-dimension quotient map of Theorem~\ref{thm:LaneA_associated_graded_identity} is fixed once and for all and survives cap enlargement.

\begin{lemma}[Canonical block data are cap-compatible]\label{lem:LaneA_block_cap_compat}
Fix an exact-dimension symmetry block $\mathscr V_{\Delta,\epsilon,\lambda}$.
Then the canonical local basis, canonical flowed basis, and canonical block norm from Definition~\ref{def:LaneA_canonical_block_norm}
are literally the same whether the block is viewed inside the truncation $\Delta\le \Delta_{\max}$ or inside any larger truncation $\Delta\le \Delta'_{\max}$.
In particular, operator bounds recorded in the canonical block norm are statements about the block itself and not about the ambient engineering cap.
\end{lemma}

\begin{proof}
The canonical local basis is fixed block-by-block before any truncation is chosen.
The canonical flowed basis is obtained from that same list of contraction formulas by replacing each curvature/derivative jet with its flowed counterpart, so it is equally independent of the ambient cap.
Finally, Lemma~\ref{lem:coherent_Z} says that the coherent renormalization scheme preserves already normalized blocks when the cap is enlarged.
Hence all blockwise coefficient and operator-norm statements are compatible under cap enlargement.
\end{proof}

\begin{theorem}[Coefficient extraction theorem for arbitrary exact-dimension symmetry blocks]\label{thm:LaneA_block_coefficient_theorem}
Fix an exact-dimension symmetry block $\mathscr V_{\Delta,\epsilon,\lambda}$ and choose the canonical local and flowed bases as above.
Let
\[
C^{\mathrm{UV,tree}}_{\Delta,\epsilon,\lambda}
\]
be the tree-level exact-dimension diagonal block of the UV matching matrix obtained by applying the coefficient-extraction map of
Lemma~\ref{lem:LaneA_block_coeff_extract} to the small-flow images of the canonical basis classes.
Then
\[
C^{\mathrm{UV,tree}}_{\Delta,\epsilon,\lambda}=I.
\]
In particular, the exact-dimension diagonal UV matching problem on this arbitrary block is an intrinsic finite-dimensional linear-algebra problem attached to the quotient block itself,
and the same matrix is seen in every engineering cap containing that block.
\end{theorem}

\begin{proof}
Apply Lemma~\ref{lem:LaneA_block_coeff_extract} to each canonical basis vector $e_q$.
Its small-flow image has the same exact-dimension coefficient vector as $e_q$ itself, namely the $q$th standard basis vector of $\mathbb R^m$.
Hence the extracted tree-level matching matrix has columns $e_1^{\mathrm{std}},\dots,e_m^{\mathrm{std}}$ and is therefore the identity.
Lemma~\ref{lem:LaneA_block_cap_compat} shows that the canonical bases and norms are unchanged under cap enlargement,
so the same matrix is obtained in every engineering cap containing the block.
\end{proof}


\begin{lemma}[Reference scalar normalization for the non-triviality witness]\label{lem:LaneA_scalar_reference_norm}
Fix the canonical CP-even scalar exact-dimension block and choose a nonzero positive-time test function
\[
\phi_{\ast,S}\in C_c^{\infty}(\{x_0>0\})
\]
from the reference family used to normalize that block.
The coherent Lane~A normalization may be fixed so that for every dyadic scale $t_j$ and every engineering cap containing the scalar block,
if $X^{S}_j(\phi_{\ast,S})$ denotes the corresponding centered renormalized scalar approximant, then
\begin{equation}\label{eq:LaneA_scalar_ref_norm}
\omega\big(\Theta X^{S}_j(\phi_{\ast,S})\,X^{S}_j(\phi_{\ast,S})\big)=1.
\end{equation}
This normalization is compatible under cap enlargement.
\end{lemma}

\begin{proof}
For the CP-even scalar block the reference Gram matrix from the normalization scheme of
Proposition~\ref{prop:flow_operator_multiplet_cauchy} and Theorem~\ref{thm:field_algebra_truncation}
is $1\times 1$.
Nondegeneracy of that Gram matrix therefore says that for some nonzero positive-time test function its single entry is a strictly positive scalar.
Because the block is one dimensional, the remaining freedom is multiplication by a nonzero scalar.
Rescale the canonical scalar basis once and for all so that the chosen reference value equals $1$.
The resulting condition is exactly \eqref{eq:LaneA_scalar_ref_norm}.
Compatibility under cap enlargement is then immediate from Lemma~\ref{lem:LaneA_block_cap_compat} together with the coherent renormalization statement of Lemma~\ref{lem:coherent_Z},
which says that already normalized blocks are unchanged when a larger engineering cap is introduced.
\end{proof}

\begin{remark}[Why the sample scalar/tensor/loop sectors remain in the paper]
The audited CP-even scalar, CP-odd scalar, small-loop scalar, traceless-tensor, and concrete dim-$6$ loop sectors now have a narrower role.
They are explicit sample calculations and normalization checks against the canonical block protocol above.
The load-bearing widening step for Lane~A is the finite-truncation theorem in Section~\ref{sec:sharp_local}, not the existence of a long ad hoc list of individually certified sectors.
\end{remark}
